

\subsubsection*{Cross the next boundary without replaying old payments}
The following-period book now contains current and deferred assessments, dated
payments against distinct tax liabilities, and the history of a replaced
assessment. Carrying this book into another period must preserve those histories
as well as its closing balances. The earlier carry caller cannot stand in for
that operation: it consumes a different predecessor and does not represent the
separate liability vintages.

The new boundary requires the recorded current and deferred assessments to reach
the financial endpoint, with the loss cursor at the same date and the closing
book dated the preceding day. The next financial interval must be contiguous.
This establishes the timing of a provisional carry, while unresolved source,
event-coverage and accounting findings remain attached to it. A caller's
assertion that those findings are true does not establish annual readiness.

In the fresh-deferred example, recorded profit is $-17$ and retained earnings
has a debit-positive balance of $-15/4$. Closing the income accounts therefore
sets retained earnings to $-15/4-(-17)=53/4$. All income accounts, including
current and deferred tax expense, start the new period at zero. Cash stays at
$715/2$, the deferred asset and liability stay at 5 each, and current-tax payable
stays at $35/4$. Assets, liabilities and total equity are unchanged. The boundary
journal balances, and the new opening has neither profit nor cash movement.

The two tax vintages retain their complete settlement and supersession records.
Separate new-period records carry outstanding amounts of $15/4$ and 5 under the
same liability identities. Each records the earlier vintage it came from, its
new financial period and an initially empty list of current-period settlements.
Earlier payments remain in the historical records; they do not enter the new
period's cash flows. A fully paid vintage is retained with a zero opening balance
rather than discarded.

The prior flows, tax adjustments, recognition calculation, measurement and
recovery capacity are archived with the closing book. Components, inventory lots,
lease-right dates, loss quantities and lifetime event identities remain intact.
Both current and deferred tax in the new period are unassessed; carrying a
recognition schedule does not renew its capacity or findings. This boundary
profile accepts no subsequent events. A compatible continuation must first
reconcile new payments against the new-period liability records while keeping
all earlier settlements historical. General annual execution and actual fiscal
completeness remain separate obligations.
